\documentclass[12pt,a4paper]{article}

% ==================================================
% 1. CONFIGURATION
% ==================================================

% Enable / disable title page
\newif\ifshowtitlepage
\showtitlepagetrue
% \showtitlepagefalse

% Enable / disable table of contents
\newif\ifshowtoc
\showtocfalse

% Document information
\newcommand{\worktitle}{Documentation Methodology}
\newcommand{\worksubtitle}{NAO Robot Exercise System}
\newcommand{\teamname}{NAOMI}
\newcommand{\subject}{NAO Robot Exercise System}

\newcommand{\university}{Slovak University of Technology in Bratislava}
\newcommand{\faculty}{Faculty of Informatics and Information Technologies}

\newcommand{\academicyear}{2026/2027}
\newcommand{\location}{Bratislava}
\newcommand{\documentyear}{2026/2027}

% Team leaders
\newcommand{\leaderone}{Ing. Jakub Perdek, PhD.}
\newcommand{\leadertwo}{ doc. Ing. Fedor Lehocki PhD.}

% Team members
\newcommand{\memberone}{Bc. Akhmadinurov Amal}
\newcommand{\membertwo}{Bc. Danilov Denis}
\newcommand{\memberthree}{Bc. Dmitriev Anton}
\newcommand{\memberfour}{Bc. Gainetdinov Ruslan}
\newcommand{\memberfive}{Bc. Kuzin Timofei}
\newcommand{\membersix}{Bc. Tikhoblazhenko Alena}
\newcommand{\memberseven}{Bc. Zimnitski Aliaksei}


% ==================================================
% 2. PACKAGES
% ==================================================

\usepackage[
    a4paper,
    top=3cm,
    bottom=2.7cm,
    left=2.8cm,
    right=2.8cm,
    headheight=45pt,
    headsep=20pt
]{geometry}

\usepackage[T1]{fontenc}
\usepackage[utf8]{inputenc}
\usepackage[english]{babel}
\usepackage{lmodern}

\usepackage{graphicx}
\graphicspath{{../assets/}{assets/}}
\usepackage[table]{xcolor}
\usepackage{fancyhdr}
\usepackage{titlesec}
\usepackage{setspace}
\usepackage{lastpage}
\usepackage{array}
\usepackage{float}
\usepackage{caption}
\usepackage{booktabs}
\usepackage{tabularx}
\usepackage{lipsum}
\usepackage{csquotes}

% Hyperlinks
\usepackage[hidelinks]{hyperref}


% ==================================================
% 3. COLORS
% ==================================================

\definecolor{primary}{HTML}{000000}
\definecolor{secondary}{HTML}{000000}
\definecolor{muted}{HTML}{000000}
\definecolor{linecolor}{HTML}{000000}
\definecolor{naomiTableHeader}{HTML}{17445A}
\definecolor{naomiTableStripe}{HTML}{EDF7FA}
\definecolor{naomiTableRule}{HTML}{66BBD0}

% Ragged-right columns keep narrow table cells readable.
\newcolumntype{L}[1]{>{\raggedright\arraybackslash}p{#1}}
\newcolumntype{Y}{>{\raggedright\arraybackslash}X}


% ==================================================
% 4. TYPOGRAPHY
% ==================================================

\setstretch{1.2}

\setlength{\parindent}{0pt}
\setlength{\parskip}{7pt}

% Section
\titleformat{\section}
    {\Large\bfseries\color{primary}}
    {\thesection}
    {0.8em}
    {}

% Subsection
\titleformat{\subsection}
    {\large\bfseries\color{primary}}
    {\thesubsection}
    {0.8em}
    {}

% Subsubsection
\titleformat{\subsubsection}
    {\normalsize\bfseries\color{primary}}
    {\thesubsubsection}
    {0.8em}
    {}

\titlespacing*{\section}{0pt}{20pt}{10pt}
\titlespacing*{\subsection}{0pt}{14pt}{7pt}
\titlespacing*{\subsubsection}{0pt}{12pt}{5pt}

% Captions
\captionsetup{
    font=small,
    labelfont=bf,
    justification=centering
}


% ==================================================
% 5. HEADER AND FOOTER
% ==================================================

\pagestyle{fancy}
\fancyhf{}

% Header logo
\fancyhead[L]{
    \includegraphics[
        width=4cm,
        height=1.15cm,
        keepaspectratio
    ]{horizontal.png}
}

% Header information
\fancyhead[R]{
    \begin{tabular}{r}
        {\footnotesize\bfseries \worksubtitle}\\[3pt]
        {\scriptsize Academic Year \academicyear}
    \end{tabular}
}

% Header line
\renewcommand{\headrulewidth}{0.7pt}

% Footer left
\fancyfoot[L]{
    {\scriptsize
    \teamname\enspace|\enspace\subject}
}

% Footer right
\fancyfoot[R]{
    {\scriptsize
    \thepage\enspace/\enspace\pageref{LastPage}}
}

% Footer line
\renewcommand{\footrulewidth}{0.4pt}

% Plain pages
\fancypagestyle{plain}{
    \fancyhf{}

    \fancyhead[L]{
        \includegraphics[
            width=4cm,
            height=1.15cm,
            keepaspectratio
        ]{horizontal.png}
    }

    \fancyhead[R]{
        \begin{tabular}{r}
            {\footnotesize\bfseries \worksubtitle}\\[3pt]
            {\scriptsize Academic Year \academicyear}
        \end{tabular}
    }

    \fancyfoot[L]{
        {\scriptsize
        \teamname\enspace|\enspace\subject}
    }

    \fancyfoot[R]{
        {\scriptsize
        \thepage\enspace/\enspace\pageref{LastPage}}
    }

    \renewcommand{\headrulewidth}{0.7pt}
    \renewcommand{\footrulewidth}{0.4pt}
}




% ==================================================
% 6. TITLE PAGE
% ==================================================

\newcommand{\makemytitlepage}{
    \begin{titlepage}
        \thispagestyle{empty}
        \centering

        % ==========================================
        % UNIVERSITY LOGO
        % ==========================================

        \vspace*{0.3cm}

        \includegraphics[
            width=6cm,
            height=5cm,
            keepaspectratio
        ]{vertical.png}

        \vspace{0.5cm}

        {\large\bfseries
        \university\par}

        \vspace{0.2cm}

        {\normalsize
        \faculty\par}

        % ==========================================
        % PROJECT TITLE
        % ==========================================

        \vspace{1.7cm}

        {\small\MakeUppercase{\worksubtitle}\par}

        \vspace{0.5cm}

        \begin{minipage}{0.95\textwidth}
            \centering

            {\fontsize{25}{32}\selectfont\bfseries
            \worktitle\par}
        \end{minipage}

        % ==========================================
        % TEAM INFORMATION
        % ==========================================

        \vfill

        \begin{minipage}{0.9\textwidth}
            \centering
            \begin{minipage}[t]{0.46\linewidth}
                \centering
                {\normalsize\bfseries Team Leaders\par}
                \vspace{0.25cm}
                {\small
                \leaderone\par
                \leadertwo\par}
            \end{minipage}%
            \hfill
            \begin{minipage}[t]{0.48\linewidth}
                \centering
                {\normalsize\bfseries Team Members\par}
                \vspace{0.25cm}
                {\small
                \memberone\par
                \membertwo\par
                \memberthree\par
                \memberfour\par
                \memberfive\par
                \membersix\par
                \memberseven\par}
            \end{minipage}
        \end{minipage}

        % ==========================================
        % BOTTOM INFORMATION
        % ==========================================

        \vfill

        {\small
        \location\enspace\textbullet\enspace\academicyear\par}

    \end{titlepage}
}

% ==================================================
% 7. IMAGE PLACEHOLDER
% ==================================================

% Usage:
% \imageplaceholder{Text}{Caption}

\newcommand{\imageplaceholder}[2]{

    \begin{figure}[H]
        \centering

        \begingroup
        \setlength{\fboxsep}{0pt}

        \fbox{
            \begin{minipage}[c][5cm][c]{0.85\textwidth}

                \centering

                {\Large\bfseries IMAGE}

                \vspace{0.3cm}

                {\normalsize #1}

            \end{minipage}
        }

        \endgroup

        \caption{#2}

    \end{figure}
}


% ==================================================
% 8. REAL IMAGE COMMAND
% ==================================================

% Usage:
% \documentimage{path}{caption}{label}

\newcommand{\documentimage}[3]{

    \begin{figure}[htbp]
        \centering

        \includegraphics[
            width=0.85\textwidth,
            keepaspectratio
        ]{#1}

        \caption{#2}
        \label{#3}

    \end{figure}
}


% ==================================================
% 9. TABLES
% ==================================================

% Usage:
% \begin{naomitable}{Caption}{label}
%     \begin{tabularx}{\textwidth}{L{2.8cm} Y L{2.6cm}}
%         \rowcolor{naomiTableHeader}
%         \textcolor{white}{\bfseries Name} & ... \\
%         ...
%         \arrayrulecolor{naomiTableRule}\bottomrule[1pt]
%     \end{tabularx}
% \end{naomitable}

\newenvironment{naomitable}[2]{%
    \begin{table}[htbp]
        \centering
        \caption{#1}
        \label{#2}
        \begingroup
        \small
        \renewcommand{\arraystretch}{1.4}
        \setlength{\tabcolsep}{7pt}
        \rowcolors{2}{naomiTableStripe}{white}
}{%
        \endgroup
    \end{table}
}

% Usage:
% \tableplaceholder{Caption}{label}

\newcommand{\tableplaceholder}[2]{
    \begin{naomitable}{#1}{#2}
        \begin{tabularx}{\textwidth}{L{2.8cm} Y L{2.6cm}}
            \rowcolor{naomiTableHeader}
            \textcolor{white}{\bfseries Name} &
            \textcolor{white}{\bfseries Description} &
            \textcolor{white}{\bfseries Status} \\
            Item A & Sample description & Completed \\
            Item B & Sample description & In Progress \\
            Item C & Sample description & Planned \\
            \arrayrulecolor{naomiTableRule}\bottomrule[1pt]
        \end{tabularx}
    \end{naomitable}
}


% ==================================================
% 10. DOCUMENT
% ==================================================

\begin{document}

\ifshowtitlepage
    \makemytitlepage
\fi

\ifshowtoc
    \tableofcontents
    \clearpage
\fi

% ==================================================
% 1. PURPOSE OF THE DOCUMENT
% ==================================================

\section{Purpose of the Document}

This document defines how the NAOMI team creates, reviews, stores, and
maintains project documentation. Its purpose is to keep technical information,
meeting outcomes, and sprint records understandable and accessible to all
team members throughout the development of the team project.

Documentation should explain not only what was implemented, but also important
decisions, known limitations, and instructions needed to continue the work.
It should be concise enough to remain useful and detailed enough that another
team member can understand the documented activity without additional
explanations from its author.

% ==================================================
% 2. SCOPE AND RESPONSIBILITIES
% ==================================================

\section{Scope and Responsibilities}

The methodology applies to documentation produced within the NAOMI team
project, including technical documentation, meeting notes, sprint
retrospectives, internal methodologies, and supporting guides.

The team member responsible for a task is normally responsible for keeping
its related documentation accurate. Documentation may also be prepared by
another member, provided that the person familiar with the implementation
reviews the technical details. All team members are encouraged to report
missing or outdated information.

The documentation is intended primarily for the team, project supervisors,
and other readers who need to understand the project's progress and results.

% ==================================================
% 3. TECHNICAL DOCUMENTATION
% ==================================================

\section{Technical Documentation}

Technical documentation describes the parts of the system developed by the
team, their purpose, and how they work together. Depending on the task,
it may cover system architecture, interfaces, data structures, APIs,
configuration, installation, or testing procedures.

When a feature is implemented or significantly changed, the responsible team
member should update the relevant documentation. The description should
include the goal of the feature, its main components, dependencies, and any
important assumptions or restrictions. Diagrams, screenshots, and short code
examples can be used when they make the explanation easier to follow.

Before the documentation is considered complete, the developer familiar with
the corresponding code should check it for technical accuracy. Relevant
documentation should also be considered during code review so that the
written description does not contradict the implementation.

Where the project contains several independent components, each component
may have its own documentation. A higher-level overview should explain how
these components fit together.


% ==================================================
% 4. MEETING RECORDS
% ==================================================

\section{Meeting Records}

During each team meeting, one team member is
responsible for taking written notes. 

The meeting notes should include the main topics
discussed, identified problems, important decisions,
and tasks assigned to individual team members.

Each meeting record should contain:
\begin{itemize}
    \item Date of the meeting and list of participants.
    \item Main topics discussed during the meeting.
    \item Decisions made by the team.
    \item Identified problems or potential risks.
    \item Action items, responsible members,
          and expected deadlines.
\end{itemize}

After the meeting, the notes are made available
to all team members. Any new tasks resulting
from the discussion should be added to Jira
or linked to existing issues.

Meeting notes are stored with the team's
project documentation so that members can
review previous discussions and decisions
when needed.

% ==================================================
% 5. INTERNAL METHODOLOGIES
% ==================================================

\section{Internal Methodologies}

In addition to technical descriptions, the NAOMI team maintains short
methodology documents for recurring activities. These documents help team
members follow a consistent approach without introducing unnecessary
administrative work.

Examples include backlog management in Jira, the Scrumban workflow,
code review, testing, version control, and documentation itself.
Each methodology should clearly describe its purpose, scope, main steps,
and responsibilities. It should be revised when the team's working
practices change.

The current version of a methodology should be easy to find in the team's
shared documentation and, where appropriate, linked from the team website.

% ==================================================
% 6. SPRINT RETROSPECTIVE DOCUMENTATION
% ==================================================

\section{Sprint Retrospective Documentation}

At the end of a sprint, the team reflects on its work and discusses possible
improvements. Although NAOMI uses Scrumban, regular retrospectives help the
team evaluate both sprint outcomes and the flow of tasks across the Jira
board.

The discussion should cover the following questions:

\begin{itemize}
    \item What went well during the sprint?
    \item Which planned activities were not completed, and why?
    \item What problems or blockers did the team encounter?
    \item What should be improved in the next sprint?
\end{itemize}

One participant should record the most important observations and decisions.
The retrospective does not need to reproduce every statement verbatim;
it should summarise common findings and concrete improvement proposals.

\subsection{Retrospective Outcomes}

The final record should distinguish between observations and agreed actions.
Where an improvement requires follow-up work, it should have a responsible
person and, if appropriate, a related Jira issue. At the next retrospective,
the team can check whether previously agreed changes were useful.

Team members who could not attend should have access to the summary and an
opportunity to provide additional comments.

% ==================================================
% 7. OTHER DOCUMENTS AND STORAGE
% ==================================================

\section{Other Documents and Storage}

Additional materials may be created during the project, such as setup guides,
user instructions, research notes, test reports, diagrams, or demonstrations.
These materials should be stored alongside the relevant project documentation
rather than remaining only in private messages or on individual computers.

Files should use descriptive names and be organised by topic. When an updated
version replaces an older one, the team should avoid keeping several
unlabelled copies that could cause confusion. Documents referenced in Jira
should be linked from the corresponding issues whenever this is useful.

Project documentation intended for broader access should also be linked or
published through the NAOMI team website, according to the team's agreed
sharing rules. Internal notes or information that should not be public should
remain in the appropriate restricted team workspace.

The goal is to keep documentation current, searchable, and useful for future
team members and project evaluation.

\end{document}
